\chapter{DOGMA primitives and state semantics}
Calling a model DNA-native does not define what it computes. Which input
symbols select operations? What persists between positions? Does padding
advance the state? What happens when a sequence is read from the opposite
strand? These questions must have executable answers before an architecture
can support meaningful training or benchmarking claims.

We use the book's taxonomy: DOGMA names the non-Transformer, DNA-native
research direction and engine; HermonDNA names the Transformer-based direction;
Evolutor is the broader framework. This chapter proposes one small state
reference for DOGMA research. It is neither a completed DOGMA architecture
nor a trained genomic model. Its deliberately limited representation makes
its invariants and information loss visible.

\section{Give each adjacent pair a stable address}
Let $\Sigma=\{A,C,G,T\}$ and order the 16 dinucleotides lexicographically
using that base order. A memory vector $h\in\mathbb R^{16}$ has one channel
per pair. Maintain a previous symbol $p$ and valid-position count $t$.
The initial state is $(0,\varepsilon,0)$.

For a valid incoming base $x$, use
\begin{equation}
 h'=\lambda h+1_{p\ne\varepsilon}e_{px},\qquad
 p'=x,\qquad t'=t+1,\quad 0\leq\lambda\leq1,
 \label{eq:transition}
\end{equation}
where $e_{px}$ is the one-hot vector selecting the ordered-pair channel.
For the first base there is no previous symbol, so nothing is added.
The pair selects a memory address, not a biochemical binding reaction.
The variable $\lambda$ is a retention coefficient, not a measured reaction rate.

\Plate{01-loci}{Previous and current bases select one of 16 ordered-pair
channels. For the transition G followed by A, only GA receives a unit increment
after all channels are multiplied by retention. Addressing is directional:
GA and AG are distinct channels.}{fig:loci}

\Listing{State}
The frozen dataclass and tuple memory prevent later steps from overwriting
an earlier state. Immutability makes replay comparisons easier; it is not
a substitute for validation. The companion rejects negative or nonfinite
memory, invalid previous symbols, and inconsistent zero-position state.
The state validator checks representation validity, not whether every supplied
state is reachable from the initial condition under one fixed $\lambda$.

\section{Implement the transition before adding learned parameters}
\Listing{step}
Every event contains the before state, consumed symbol, selected pair and
after state. A trace can therefore be replayed by applying the same transition
and checking equality, rather than trusting a textual explanation attached
afterward. Padding emits no event and returns the identical state object.
An unknown unmasked base is rejected instead of silently converted to A or
treated as an absence of data.

For ACGAC with $\lambda=1$, the generated trace is
\begin{center}\input{\Generated/trace.tex}\end{center}
The previous column records the symbol retained \emph{after} each step.
The selected pairs are AC, CG, GA and AC. Hence final counts are AC twice,
CG once and GA once. Five bases create four adjacent pairs, not five.

The complete scanner is similarly explicit:
\Listing{scan}
It returns both final state and immutable event sequence. Passing a final
state into the next call continues the same record. Calling with no state
begins a new record. The function validates retention even on empty input,
so invalid configuration cannot be hidden by an empty batch.

\section{Chunking, padding and reset are different operations}
For a fixed $\lambda$, sequential composition gives
\begin{equation}
 \operatorname{scan}(uv,s)=
 \operatorname{scan}\bigl(v,\operatorname{scan}(u,s)\bigr),
 \label{eq:chunk}
\end{equation}
where the inner result denotes its final state. Concatenating the two event
sequences also reproduces the uninterrupted trace. This property follows
because each step depends only on its current state and symbol, not the
call boundary. Changing retention between chunks changes the computation.

\Plate{03-step}{A valid base updates memory and previous symbol. A padding
slot is the identity on all state fields and emits no event. A new biological
record instead starts from an explicit fresh state. Skipping an unknown base
as padding can create a pair across a real sequence gap.}{fig:step}

The mask denotes storage padding, not an uncertain biological nucleotide.
For symbols A, N, C with mask true, false, true, the model counts AC because
the middle slot is absent from the logical sequence. If N represents an
unknown base in a real record, this is usually the wrong semantics: the
adjacency of A and C is not known. Alternatives include a separate unknown
symbol, an explicit boundary, a distribution over bases or rejection. Each
requires a new contract and tests.

Similarly, concatenating two independent records without resetting creates a
spurious pair across their boundary. A reset clears memory, previous symbol
and position. It is not equivalent to zero retention: at $\lambda=0$, the
previous symbol still creates the next pair. A batching interface should
carry record identity separately from a padding mask.

\section{Separate chemical polarity from string orientation}
A DNA duplex is antiparallel. If one strand is written $5'$ to $3'$ as
ACGAC, the aligned complementary strand runs $3'$ to $5'$ as TGCTG.
Reading that complementary strand in its own $5'$ to $3'$ direction gives
GTCGT. This last operation is the \Term{reverse complement}, not just a
characterwise complement.

\Plate{02-strands}{The upper and aligned lower strands have opposite chemical
polarity. The software reverse-complement string is obtained by reading the
lower strand from its 5-prime end. Pair channels must reverse order and
complement both bases; complementing without reversal implements a different
transformation.}{fig:strands}
\Listing{reverse_complement}

Let $R$ denote reverse complementation of a sequence and $P$ its channel
permutation: pair $ab$ maps to $c(b)c(a)$, where $c$ complements a base.
For example, AC maps to GT, GA maps to TC, and CG maps to itself.
Applying either transformation twice returns the original: $R^2=I$ and
$P^2=I$.
\Listing{permute_pairs}

\section{Prove a symmetry, then construct its failure}
With $\lambda=1$, the final memory $H(s)$ is a histogram of adjacent pairs.
Every adjacent pair in $s$ corresponds to exactly one reverse-complemented
pair in $R(s)$. This bijection proves
\begin{equation}
 H(R(s))=P H(s).
 \label{eq:equiv}
\end{equation}
This is \Term{equivariance}: features change by a known permutation when
the input orientation changes. It is not invariance, which would require
the feature vector itself to stay unchanged.

\begin{center}\input{\Generated/counts.tex}\end{center}
The generated table maps the reverse-complement memory back to forward
channel coordinates before comparing it. Equality without that mapping
would compare different biological orientations in different feature bases.

\Plate{04-equivariance}{For unit retention, each observed adjacent pair
contributes equally and reverse complementation permutes the histogram.
For AAC at retention one half, older AA receives weight one half and newer
AC receives weight one. Reversing orientation reverses their ages, breaking
the one-pass symmetry.}{fig:equivariance}

For nonunit retention, earlier pairs receive more decay. AAC at $\lambda=0.5$
has $H_{AA}=0.5,H_{AC}=1$. Its reverse complement GTT has
$H_{GT}=0.5,H_{TT}=1$. Mapping back yields AA equal to one and AC equal
to one half, not the forward memory. The addressing rule respects complement
mapping, but the direction-dependent weighting does not. A model can fail
strand symmetry even when every local base transformation looks correct.

One exact full-record repair is group averaging:
\begin{equation}
 Z(s)=\tfrac12\bigl(H(s)+P H(R(s))\bigr).
\end{equation}
Because $P^2=I$, we have $Z(R(s))=PZ(s)$ for any $H$ with these channel
dimensions. This proof does not require $H$ itself to be equivariant.
\Listing{bidirectional}

The repair reads the full record in both orientations. It is not an online
causal update and should not be sold as one. A scalar output $w^\top Z$
is invariant only if its head satisfies the corresponding symmetry, for
example $w=P w$. Strand-specific tasks may instead require equivariant
outputs that swap labels. The correct property depends on what the labels
mean, not on a universal preference for invariance.

\section{Memory mass exposes retention and information loss}
Let $M_n$ be the sum of all channels after $n$ valid bases from a fresh state.
Then $M_0=M_1=0$ and $M_n=\lambda M_{n-1}+1$ for $n\geq2$.
Consequently,
\begin{equation}
 M_n=\sum_{j=0}^{n-2}\lambda^j=
 \begin{cases}
 n-1,&\lambda=1,\ n\geq1,\\
 (1-\lambda^{n-1})/(1-\lambda),&0\leq\lambda<1,\ n\geq1.
 \end{cases}
\end{equation}
The empty case is defined separately as zero. For $\lambda<1$, mass is
bounded by $1/(1-\lambda)$, while unit retention accumulates without bound
in the exact arithmetic model.

\Plate{05-forgetting}{Generated memory mass for repeated A inputs at three
retention settings. The sequence identity does not affect total mass in this
reference; every adjacent pair contributes one before future decay. Bounded
mass does not imply that all useful sequence information is preserved.}{fig:forgetting}

At unit retention, AACA and ACAA have the same final pair counts (AA, AC,
CA once each), the same previous base A and the same position four.
The complete final states are identical although the sequences differ.
No readout of that state can distinguish them. This is a constructive
information-loss proof, not a failure of training. Learning an output head
cannot recover information that the state has already discarded.

\section{Resource and numerical claims need their own tests}
The reference updates 16 channels per valid step: $O(16n)$ scalar work for
$n$ bases and 16 memory slots in the recurrent vector. Materialized events
retain before and after states, so a full trace requires $O(16n)$ stored
scalar entries up to sharing and implementation overhead. The position counter
also grows: its exact bit length is $O(\log n)$. Fixed vector width is not
fixed exact information capacity at arbitrary sequence length.

Python floats cannot retain every unit increment once a channel exceeds
the consecutive-integer range of binary64. Fractional retention introduces
rounding much earlier, so tests use declared tolerances where needed.
An accelerated implementation must compare chunking, masking, reset and
orientation behavior against the reference, not merely match one aggregate
loss. Vectorizing across records must never leak previous symbols between them.

The tests exhaustively compare pair counts with an independent counter for
short sequences, test reverse-complement involution, check all chunk splits
at several retentions, verify memory mass, and include the explicit decay
counterexample. They establish the toy operator's contract, not genomic
prediction quality, hardware throughput or learned biological mechanisms.

\section{Connect to research without borrowing its claims}
Caduceus distinguishes bidirectional sequence processing from reverse-complement
equivariance and constructs symmetry-aware DNA sequence models \cite{caduceus}.
That architectural distinction motivates the questions asked here. Our pair
histogram and two-pass average are not a Caduceus implementation and carry
none of its experimental performance claims. A useful next DOGMA experiment
would compare precisely specified state operators under matched data splits,
parameter budgets and label transformations before making a performance claim.

The primitive's value is its inspectability. It gives a baseline with exact
pair semantics, explicit loss of order, known symmetry conditions and a replay
contract. More expressive operators can now be justified by showing which
limitation they remove and what new cost or verification burden they add.

\section{Exercises with worked answers}
\begin{enumerate}
\item Scan ACG at unit retention. \textbf{Answer:} AC and CG each equal
one; previous is G and position is three. All other channels are zero.
\item Why is AG not selected by G then A? \textbf{Answer:} Addresses are
ordered previous--current pairs, so the selected channel is GA.
\item Split ACG after A. \textbf{Answer:} Carry previous A into the
second chunk; consuming C must count AC. Resetting would lose that pair.
\item Pad between A and C. \textbf{Answer:} A masked slot changes
nothing; C still creates AC. This assumes the slot is not a biological base.
\item Why is unknown N not ordinary padding? \textbf{Answer:} A real
unknown position interrupts known adjacency. Skipping it asserts a different
logical sequence.
\item Show retention zero is not reset. \textbf{Answer:} After A, a
valid C with zero retention still adds AC using the retained previous base.
\item Reverse-complement ACGAC. \textbf{Answer:} Complement gives
TGCTG in the aligned opposite polarity; reversing gives GTCGT.
\item Map channel GA under reverse complementation. \textbf{Answer:}
Reverse AG and complement to TC. Applying the mapping again gives GA.
\item Explain decay's symmetry failure. \textbf{Answer:} Reverse orientation
changes the age of each pair, so fixed forward decay assigns different weights
even after channel permutation.
\item Prove the averaged feature is equivariant. \textbf{Answer:}
$Z(Rs)=(H(Rs)+PH(s))/2=P(H(s)+PH(Rs))/2=PZ(s)$.
\item Compute mass for four bases at retention one half.
\textbf{Answer:} $1+1/2+1/4=1.75$, independent of which canonical bases occur.
\item Can an output head distinguish AACA from ACAA at unit retention?
\textbf{Answer:} Not from this final state: the counts, final base and
length all agree. A richer representation or access to the trace is required.
\end{enumerate}

\section{Vocabulary and the next question}
\Term{DNA-addressed state}: software memory whose update address is selected
by declared sequence symbols. \Term{Padding identity}: a masked slot leaves
all logical state unchanged. \Term{Record reset}: explicit replacement with
an initial state at a biological record boundary. \Term{Bidirectionality}:
access to sequence context in both directions, distinct from symmetry under
reverse complementation.

The next chapter can add richer context and trainable operators while keeping
these contracts visible. Every extra capability should arrive with a worked
counterexample it resolves, a reference implementation and a test capable of
revealing the new failure modes.
